\documentclass{beamer}
\mode<presentation>
\usepackage{amsmath,amssymb,mathtools}
\usepackage{textcomp}
\usepackage{gensymb}
\usepackage{adjustbox}
\usepackage{subcaption}
\usepackage{enumitem}
\usepackage{multicol}
\usepackage{listings}
\usepackage{url}
\usepackage{graphicx} % <-- needed for images
\usepackage{xcolor}
\def\UrlBreaks{\do\/\do-}

\usetheme{Boadilla}
\usecolortheme{lily}
\setbeamertemplate{footline}{
	\leavevmode%
	\hbox{%
		\begin{beamercolorbox}[wd=\paperwidth,ht=2ex,dp=1ex,right]{author in head/foot}%
			\insertframenumber{} / \inserttotalframenumber\hspace*{2ex}
	\end{beamercolorbox}}%
	\vskip0pt%
}
\setbeamertemplate{navigation symbols}{}

\lstset{
	frame=single,
	breaklines=true,
	columns=fullflexible,
	basicstyle=\ttfamily\tiny    % tiny font so code fits
}

\numberwithin{equation}{section}

% ---- your macros ----
\providecommand{\nCr}[2]{\,^{#1}C_{#2}}
\providecommand{\nPr}[2]{\,^{#1}P_{#2}}
\providecommand{\mbf}{\mathbf}
\providecommand{\pr}[1]{\ensuremath{\Pr\left(#1\right)}}
\providecommand{\qfunc}[1]{\ensuremath{Q\left(#1\right)}}
\providecommand{\sbrak}[1]{\ensuremath{{}\left[#1\right]}}
\providecommand{\lsbrak}[1]{\ensuremath{{}\left[#1\right.}}
\providecommand{\rsbrak}[1]{\ensuremath{\left.#1\right]}}
\providecommand{\brak}[1]{\ensuremath{\left(#1\right)}}
\providecommand{\lbrak}[1]{\ensuremath{\left(#1\right.}}
\providecommand{\rbrak}[1]{\ensuremath{\left.#1\right)}}
\providecommand{\cbrak}[1]{\ensuremath{\left\{#1\right\}}}
\providecommand{\lcbrak}[1]{\ensuremath{\left\{#1\right.}}
\providecommand{\rcbrak}[1]{\ensuremath{\left.#1\right\}}}
\theoremstyle{remark}
\newtheorem{rem}{Remark}
\newcommand{\sgn}{\mathop{\mathrm{sgn}}}
\providecommand{\abs}[1]{\left\vert#1\right\vert}
\providecommand{\res}[1]{\Res\displaylimits_{#1}}
\providecommand{\norm}[1]{\lVert#1\rVert}
\providecommand{\mtx}[1]{\mathbf{#1}}
\providecommand{\mean}[1]{E\left[ #1 \right]}
\providecommand{\fourier}{\overset{\mathcal{F}}{ \rightleftharpoons}}
\providecommand{\system}{\overset{\mathcal{H}}{ \longleftrightarrow}}
\providecommand{\dec}[2]{\ensuremath{\overset{#1}{\underset{#2}{\gtrless}}}}
\newcommand{\myvec}[1]{\ensuremath{\begin{pmatrix}#1\end{pmatrix}}}
\newcommand{\mydet}[1]{\ensuremath{\begin{vmatrix}#1\end{vmatrix}}}

\newenvironment{amatrix}[1]{%
	\left[\begin{array}{@{}*{#1}{c}|c@{}}
	}{%
	\end{array}\right]
}

\newcommand{\myaugvec}[2]{\ensuremath{\begin{amatrix}{#1}#2\end{amatrix}}}
\let\vec\mathbf
% ---------------------

\title{12.159}
\author{EE25BTECH11042 - Nipun Dasari}

\begin{document}
	
	\begin{frame}
		\titlepage
	\end{frame}
	
	\begin{frame}{Question}
		The determinant
		\begin{align}
			\mydet{1+b & b & 1 \\ b & 1+b & 1 \\ 1 & 2b & 1}
		\end{align}
		evaluates to
		\begin{enumerate}[label=\alph*)]
			\item $0$
			\item $2b(b-1)$
			\item $2(1-b)(1+2b)$
			\item $3b(1+b)$
		\end{enumerate}
	\end{frame}
	
	\begin{frame}{Theoretical solution}
		Let the determinant of the given matrix be $\Delta$. We will find its value by converting the matrix to Row Echelon Form (REF).
		\begin{align}
			\Delta = \mydet{1+b & b & 1 \\ b & 1+b & 1 \\ 1 & 2b & 1}
		\end{align}
		First, apply the operation $R_2 \to R_2 - R_1$.
		\begin{align}
			\Delta = \mydet{1+b & b & 1 \\ b-(1+b) & 1+b-b & 1-1 \\ 1 & 2b & 1} = \mydet{1+b & b & 1 \\ -1 & 1 & 0 \\ 1 & 2b & 1}
		\end{align}
	\end{frame}
	
	\begin{frame}{Theoretical solution}
		Next, apply the operation $R_3 \to R_3 - R_1$.
		\begin{align}
			\Delta = \mydet{1+b & b & 1 \\ -1 & 1 & 0 \\ 1-(1+b) & 2b-b & 1-1} = \mydet{1+b & b & 1 \\ -1 & 1 & 0 \\ -b & b & 0}
		\end{align}
		Here, we can observe that the third row is a scalar multiple of the second row: $R_3 = b \cdot R_2$. This implies the rows are linearly dependent.
	\end{frame}
	
	\begin{frame}{Theoretical solution}
		To complete the reduction to REF, we perform one final operation, $R_3 \to R_3 - bR_2$, to create a zero row.
		\begin{align}
			\Delta = \mydet{1+b & b & 1 \\ -1 & 1 & 0 \\ -b - b(-1) & b-b(1) & 0-b(0)} = \mydet{1+b & b & 1 \\ -1 & 1 & 0 \\ 0 & 0 & 0}
		\end{align}
		Since the matrix in Row Echelon Form has a row of all zeros, its determinant is 0.
		\begin{align}
			\Delta = 0
		\end{align}
		This corresponds to option (a).
	\end{frame}
	
\end{document}